\documentclass[11pt,a4paper]{article}

% ------------------------------------------------------------------
%  Paquetes
% ------------------------------------------------------------------
\usepackage[utf8]{inputenc}
\usepackage[T1]{fontenc}
\IfFileExists{spanish.ldf}
  {\usepackage[spanish,es-noshorthands,es-nodecimaldot]{babel}}
  {\usepackage[english]{babel}
   \addto\captionsenglish{\renewcommand{\contentsname}{Índice}}}
\IfFileExists{newpxtext.sty}
  {\usepackage{newpxtext}\usepackage{newpxmath}}
  {\usepackage{mathpazo}}
\usepackage{microtype}
\usepackage[a4paper,top=2.8cm,bottom=2.8cm,left=2.9cm,right=2.9cm,headheight=14pt]{geometry}
\usepackage{amsmath,amssymb,amsthm,mathtools,bm}
\usepackage{booktabs,array}
\usepackage{enumitem}
\usepackage{titlesec}
\usepackage{fancyhdr}
\usepackage[hidelinks]{hyperref}

\setlength{\parindent}{1.5em}
\setlength{\parskip}{2pt}
\setlist{itemsep=1pt,topsep=4pt,leftmargin=2em}
\renewcommand{\arraystretch}{1.2}
\linespread{1.05}

% ------------------------------------------------------------------
%  Títulos y encabezados
% ------------------------------------------------------------------
\titleformat{\section}{\normalfont\Large\bfseries}{\thesection}{1em}{}
\titleformat{\subsection}{\normalfont\large\bfseries}{\thesubsection}{1em}{}
\titleformat*{\subsubsection}{\normalfont\normalsize\itshape}
\titlespacing*{\section}{0pt}{3.2ex plus 1ex minus .2ex}{1.6ex plus .2ex}
\titlespacing*{\subsection}{0pt}{2.6ex plus 1ex minus .2ex}{1.1ex plus .2ex}

\pagestyle{fancy}
\fancyhf{}
\renewcommand{\sectionmark}[1]{\markboth{\thesection.\ #1}{}}
\fancyhead[L]{\small\scshape Teoría asintótica}
\fancyhead[R]{\small\itshape\leftmark}
\fancyfoot[C]{\small\thepage}
\renewcommand{\headrulewidth}{0.4pt}

% ------------------------------------------------------------------
%  Entornos
% ------------------------------------------------------------------
\theoremstyle{plain}
\newtheorem{teorema}{Teorema}[section]
\newtheorem{proposicion}[teorema]{Proposición}
\theoremstyle{definition}
\newtheorem{definicion}[teorema]{Definición}
\newtheorem*{ejemplo}{Ejemplo}
\newtheorem*{ejemplos}{Ejemplos}
\theoremstyle{remark}
\newtheorem*{intuicion}{Intuición}
\newtheorem*{observacion}{Observación}
\newtheorem*{resumen}{Resumen}
\AtBeginDocument{\renewcommand{\proofname}{Demostración}}

% ------------------------------------------------------------------
%  Notación
% ------------------------------------------------------------------
\newcommand{\E}{\mathbb{E}}
\renewcommand{\P}{\mathbb{P}}
\newcommand{\R}{\mathbb{R}}
\DeclareMathOperator{\Var}{Var}
\DeclareMathOperator{\Cov}{Cov}
\DeclareMathOperator{\Corr}{Corr}
\DeclareMathOperator{\tr}{tr}
\DeclareMathOperator*{\plim}{plim}
\newcommand{\cp}{\xrightarrow{\;p\;}}
\newcommand{\cd}{\xrightarrow{\;d\;}}
\newcommand{\cL}[1]{\xrightarrow{\;L^{#1}\;}}
\newcommand{\cas}{\xrightarrow{\;c.s.\;}}
\newcommand{\iid}{\overset{\text{iid}}{\sim}}


\begin{document}
\thispagestyle{plain}

\begin{center}
  {\small\scshape Econometría\enspace·\enspace Notas de clase}\\[4pt]
  \rule{\textwidth}{0.6pt}\\[10pt]
  {\LARGE\bfseries Teoría asintótica}\\[6pt]
  \rule{\textwidth}{0.6pt}\\[4pt]
  {\small J.\,D.\ Cárdenas Barrantes\hfill Noviembre de 2026}
\end{center}

\vspace{1em}
{\small\tableofcontents}
\newpage

% ============================================================
\section{Introducción: el problema de fondo}
% ============================================================

En econometría queremos conocer un parámetro poblacional, por ejemplo $\beta$: un número
\emph{fijo} que describe una relación en toda la población. Nunca lo observamos. Lo que
observamos es una \emph{muestra}, y con ella calculamos un estimador $\hat\beta$.

La idea que organiza todo el curso es la siguiente: si nos hubiera tocado otra muestra,
habríamos obtenido otro $\hat\beta$. Por lo tanto, \textbf{$\hat\beta$ es una variable aleatoria}
y tiene su propia distribución: la de todos los valores que tomaría a través de todas las
muestras posibles. Eso deja dos preguntas:

\begin{enumerate}
  \item \textbf{¿Le apunta al blanco?} ¿$\hat\beta$ termina cerca del $\beta$ verdadero?
  \item \textbf{¿Cuánto se equivoca?} ¿Cómo fluctúa $\hat\beta$ alrededor de $\beta$? Sin esto no
        hay errores estándar, ni intervalos de confianza, ni pruebas de hipótesis.
\end{enumerate}

El obstáculo es que, para conocer la distribución \emph{exacta} de $\hat\beta$ con un $n$
concreto, necesitaríamos conocer la distribución de la población, que es justamente lo que
no conocemos. El truco de la teoría asintótica es preguntar qué pasa cuando $n$ es muy
grande. Ahí todo se simplifica, y esa respuesta se usa como \emph{aproximación} para la
muestra que tenemos.

\paragraph{Plan de las notas.}
\begin{itemize}
  \item \textbf{Sección \ref{sec:fund}.} Fundamentos: esperanza, varianza, covarianza, la media
        muestral, el error estándar y los grados de libertad. Dan el \emph{centro} y el
        \emph{ancho} de un estimador.
  \item \textbf{Sección \ref{sec:conv}.} Nociones de convergencia: qué significa que una sucesión
        de variables aleatorias ``se acerque'' a algo.
  \item \textbf{Sección \ref{sec:asin}.} Las herramientas asintóticas: ley de grandes números,
        teorema de la función continua, Slutsky, teorema central del límite y método delta.
        Responden las dos preguntas de arriba.
\end{itemize}

% ============================================================
\section{Fundamentos de probabilidad e inferencia}\label{sec:fund}
% ============================================================

\subsection{Población, parámetro, muestra y estimador}

\begin{definicion}[Conceptos básicos]\label{def:basicos}\leavevmode
\begin{itemize}
  \item \textbf{Población:} el universo completo que nos interesa.
  \item \textbf{Parámetro:} un número fijo que describe a la población (por ejemplo, $\mu$,
        $\sigma^2$ o $\beta$). Existe, pero es desconocido.
  \item \textbf{Muestra:} las observaciones que sí vemos, $X_1,\dots,X_n$.
  \item \textbf{Estadístico:} cualquier cálculo hecho con la muestra, como $\bar X_n$.
  \item \textbf{Estimador:} un estadístico usado para aproximar un parámetro.
\end{itemize}
\end{definicion}

\begin{intuicion}
Como la muestra salió al azar, \textbf{el estimador también es aleatorio}. Toda la inferencia
estadística consiste en entender esa aleatoriedad.
\end{intuicion}

\subsection{Variables aleatorias y distribuciones}

Una \textbf{variable aleatoria} es una regla que asigna un número a cada resultado posible del
azar: $X=X(\omega)$, donde $\omega$ representa el ``estado del mundo'' (por ejemplo, qué
muestra nos tocó). Lo que la hace aleatoria \emph{no} es tener varianza positiva, sino que su
valor dependa de $\omega$. Una constante es el caso trivial: una variable aleatoria que
siempre toma el mismo valor.

\begin{definicion}[Función de distribución acumulada (CDF)]\label{def:cdf}
La CDF de $X$ es
\[
F(u)=\P(X\le u),\qquad u\in\R .
\]
Para un vector, $X\le u$ significa que cada componente es menor o igual que la
correspondiente de $u$.
\end{definicion}

La CDF acumula probabilidad de izquierda a derecha: empieza en $0$ y termina en $1$. En una
variable discreta avanza a saltos; en una continua sube de forma suave. Describe la
distribución completa. Como describir una distribución entera es mucho, se resume con dos
números: dónde está centrada (la esperanza) y qué tan dispersa es (la varianza).

\subsection{Esperanza}

\begin{definicion}[Esperanza]\label{def:esperanza}
\[
\E[Y]=\sum_y y\,\P(Y=y)\quad\text{(caso discreto)},\qquad
\E[Y]=\int y\,f(y)\,dy\quad\text{(caso continuo)}.
\]
\end{definicion}

\begin{intuicion}
Es el centro de gravedad de la distribución: el promedio al que llegaríamos si repitiéramos
el experimento infinitas veces.
\end{intuicion}

\begin{proposicion}[Linealidad de la esperanza]\label{prop:linE}
Para constantes $a,b$ y variables $X,Y$:
\[
\E[a+bY]=a+b\,\E[Y],\qquad \E[X+Y]=\E[X]+\E[Y].
\]
La segunda vale siempre, sean $X$ e $Y$ independientes o no.
\end{proposicion}

Desplazar todos los valores en $a$ desplaza el centro en $a$; multiplicarlos por $b$
multiplica el centro por $b$. En la fórmula, las constantes salen de la suma o de la integral.

\subsection{Varianza y desviación estándar}

\begin{definicion}[Varianza]\label{def:var}
Con $\mu=\E[Y]$,
\[
\Var(Y)=\E\big[(Y-\mu)^2\big].
\]
\end{definicion}

\textbf{¿Por qué al cuadrado?} Si promediamos las desviaciones tal cual, se cancelan:
$\E[Y-\mu]=\E[Y]-\mu=0$. Al elevarlas al cuadrado todas son positivas, y las grandes pesan más.

\begin{proposicion}[Fórmula de cálculo]\label{prop:varcalc}
\[
\Var(Y)=\E[Y^2]-\mu^2 .
\]
\end{proposicion}
\begin{proof}
Desarrollamos el cuadrado: $(Y-\mu)^2=Y^2-2\mu Y+\mu^2$. Por linealidad,
$\E[(Y-\mu)^2]=\E[Y^2]-2\mu\E[Y]+\mu^2=\E[Y^2]-2\mu^2+\mu^2=\E[Y^2]-\mu^2$.
\end{proof}

\begin{proposicion}[Varianza de una transformación lineal]\label{prop:varlin}
\[
\Var(a+bY)=b^2\,\Var(Y).
\]
\end{proposicion}
\begin{proof}
La media de $a+bY$ es $a+b\mu$, así que su desviación es $(a+bY)-(a+b\mu)=b(Y-\mu)$.
Entonces $\Var(a+bY)=\E[b^2(Y-\mu)^2]=b^2\Var(Y)$.
\end{proof}

\begin{intuicion}
La constante $a$ desaparece porque mover todos los valores no cambia qué tan dispersos
están. La constante $b$ sale \emph{al cuadrado} porque la varianza mide desviaciones al
cuadrado.
\end{intuicion}

\begin{definicion}[Desviación estándar]\label{def:sd}
\[
\sigma=\sqrt{\Var(Y)}.
\]
\end{definicion}

La varianza está en unidades al cuadrado (si $Y$ está en soles, la varianza está en
soles$^2$). La raíz la devuelve a las unidades de $Y$. Por eso la desviación estándar es la
que se interpreta como ``distancia típica a la media'', la que se usa para estandarizar y la
que se reporta como error estándar.

\subsection{Covarianza, correlación e independencia}

\begin{definicion}[Covarianza y correlación]\label{def:cov}
\[
\Cov(X,Y)=\E\big[(X-\mu_X)(Y-\mu_Y)\big],\qquad
\Corr(X,Y)=\frac{\Cov(X,Y)}{\sigma_X\sigma_Y}\in[-1,1].
\]
\end{definicion}

La covarianza mide si dos variables se mueven juntas: es positiva si, cuando una está sobre
su media, la otra suele estarlo también; negativa si se mueven en sentido opuesto; cero si no
hay relación \emph{lineal}. La correlación es la covarianza estandarizada.

\begin{proposicion}[Formas equivalentes de la covarianza]\label{prop:covforms}\leavevmode
\begin{enumerate}
  \item Fórmula de cálculo: $\Cov(X,Y)=\E[XY]-\E[X]\E[Y]$.
  \item Desde la correlación: $\Cov(X,Y)=\rho\,\sigma_X\sigma_Y$.
  \item Bilinealidad: $\Cov(aX+bY,Z)=a\Cov(X,Z)+b\Cov(Y,Z)$.
  \item La varianza es un caso particular: $\Var(Y)=\Cov(Y,Y)$.
\end{enumerate}
\end{proposicion}
\begin{proof}[Demostración de (1)]
$(X-\mu_X)(Y-\mu_Y)=XY-\mu_YX-\mu_XY+\mu_X\mu_Y$. Tomando esperanza:
$\E[XY]-\mu_X\mu_Y-\mu_X\mu_Y+\mu_X\mu_Y=\E[XY]-\mu_X\mu_Y$.
\end{proof}

\begin{proposicion}[Varianza de una suma]\label{prop:varsuma}
\[
\Var(X+Y)=\Var(X)+\Var(Y)+2\Cov(X,Y),
\]
y en general
\[
\Var\Big(\sum_{i=1}^n Y_i\Big)=\sum_{i=1}^n\Var(Y_i)+\sum_{i\neq j}\Cov(Y_i,Y_j).
\]
\end{proposicion}
\begin{proof}
La desviación de la suma es la suma de las desviaciones. Al elevar al cuadrado,
$(a+b)^2=a^2+b^2+2ab$: los cuadrados dan las varianzas y los términos cruzados, las
covarianzas.
\end{proof}

\begin{proposicion}[Independencia]\label{prop:indep}
Si $X$ e $Y$ son independientes, entonces $\Cov(X,Y)=0$. Lo contrario no es cierto en general:
covarianza cero no implica independencia.
\end{proposicion}

\begin{proposicion}[Ley de varianza total]\label{prop:vartotal}
\[
\Var(Y)=\E\big[\Var(Y\mid X)\big]+\Var\big(\E[Y\mid X]\big).
\]
La variación total es la variación \emph{dentro} de los grupos más la variación \emph{entre}
los promedios de los grupos.
\end{proposicion}

\begin{definicion}[Matriz de varianzas y covarianzas]\label{def:sigma}
Para un vector aleatorio $X\in\R^k$ con media $\mu$,
\[
\Sigma=\E\big[(X-\mu)(X-\mu)'\big]=\E[XX']-\mu\mu'.
\]
La diagonal contiene las varianzas; fuera de la diagonal, las covarianzas.
\end{definicion}

\subsection{La muestra i.i.d.}

Decimos que $X_1,\dots,X_n$ son \textbf{i.i.d.} si son:
\begin{itemize}
  \item \textbf{independientes:} una observación no afecta a otra (todas las covarianzas son cero);
  \item \textbf{idénticamente distribuidas:} todas vienen de la misma población, con media $\mu$
        y varianza $\sigma^2$.
\end{itemize}

\subsection{La media muestral}\label{sec:media}

\[
\bar X_n=\frac1n\sum_{i=1}^n X_i .
\]

\begin{teorema}[Media y varianza de la media muestral]\label{teo:mediamuestral}
Si $X_1,\dots,X_n$ son i.i.d. con media $\mu$ y varianza $\sigma^2$,
\[
\E[\bar X_n]=\mu,\qquad \Var(\bar X_n)=\frac{\sigma^2}{n},\qquad
\text{desv. est.}(\bar X_n)=\frac{\sigma}{\sqrt n}.
\]
\end{teorema}
\begin{proof}
\textbf{Esperanza.} Por linealidad,
$\E[\bar X_n]=\frac1n\sum\E[X_i]=\frac1n\cdot n\mu=\mu$.

\textbf{Varianza.}
\begin{enumerate}
  \item La constante $\tfrac1n$ sale al cuadrado:
        $\Var(\bar X_n)=\frac1{n^2}\Var\big(\sum X_i\big)$.
  \item Por independencia, las covarianzas son cero; por distribución idéntica, cada varianza
        es $\sigma^2$: $\Var\big(\sum X_i\big)=n\sigma^2$.
  \item Juntando: $\Var(\bar X_n)=\frac{n\sigma^2}{n^2}=\frac{\sigma^2}{n}$, y la desviación
        estándar es su raíz, $\sigma/\sqrt n$.
\end{enumerate}
\end{proof}

\begin{intuicion}
La suma acumula incertidumbre: su varianza crece como $n$, pero su desviación estándar solo
como $\sqrt n$, porque los errores independientes \emph{se cancelan parcialmente} (unos van
hacia arriba, otros hacia abajo). Al dividir entre $n$, la incertidumbre del promedio cae como
$1/\sqrt n$. Consecuencia práctica: para reducir el error a la mitad se necesitan
\emph{cuatro} veces más datos.
\end{intuicion}

\begin{observacion}
La fórmula $\sigma^2/n$ depende del supuesto i.i.d. Si las observaciones están correlacionadas
(como en series de tiempo), las covarianzas no se anulan y hay que sumarlas: aparece la
\emph{varianza de largo plazo} y los errores estándar HAC.
\end{observacion}

\subsection{Distribución muestral y error estándar}

\begin{definicion}[Distribución muestral y error estándar]\label{def:se}
La \textbf{distribución muestral} de un estimador es su distribución a través de todas las
muestras posibles. El \textbf{error estándar} es su desviación estándar. Para la media
muestral es $\sigma/\sqrt n$, y en la práctica se estima con $s/\sqrt n$.
\end{definicion}

\begin{center}
\begin{tabular}{@{}lll@{}}
\toprule
 & Qué mide & ¿Depende de $n$? \\
\midrule
$\sigma$ & Dispersión de \emph{los datos} & No \\
$\sigma/\sqrt n$ & Dispersión \emph{del estimador} & Sí: se achica \\
\bottomrule
\end{tabular}
\end{center}

\subsection{La varianza muestral y el \texorpdfstring{$n-1$}{n-1}}

\[
s^2=\frac{1}{n-1}\sum_{i=1}^n(X_i-\bar X_n)^2 .
\]

\begin{teorema}[Insesgadez de $s^2$]\label{teo:s2}
Si las $X_i$ son i.i.d. con varianza $\sigma^2$, entonces $\E[s^2]=\sigma^2$.
\end{teorema}
\begin{proof}
\begin{enumerate}
  \item Sumando y restando $\mu$ dentro del paréntesis y desarrollando, los términos cruzados se
        simplifican y queda
        $\sum(X_i-\bar X_n)^2=\sum(X_i-\mu)^2-n(\bar X_n-\mu)^2$.
  \item $\E\big[\sum(X_i-\mu)^2\big]=n\sigma^2$ y
        $\E\big[n(\bar X_n-\mu)^2\big]=n\Var(\bar X_n)=\sigma^2$.
  \item Restando: $\E\big[\sum(X_i-\bar X_n)^2\big]=(n-1)\sigma^2$. Dividir entre $n-1$ da un
        estimador insesgado.
\end{enumerate}
\end{proof}

\begin{intuicion}[grados de libertad]
$\bar X_n$ se calcula con los mismos datos, así que está más cerca de ellos que el verdadero
$\mu$; medir las desviaciones respecto de $\bar X_n$ las subestima. Además, las desviaciones
cumplen una restricción, $\sum(X_i-\bar X_n)=0$: si conocemos $n-1$ de ellas, la última queda
determinada. Solo $n-1$ son libres. Estimar la media ``gasta'' un \textbf{grado de libertad}.

En asintótica da igual dividir entre $n$ o entre $n-1$: como $\frac{n-1}{n}\to1$, ambos
estimadores son consistentes.
\end{intuicion}

\subsection{Errores, residuos y grados de libertad en la regresión}\label{sec:gl}

En el modelo $y_i=x_i'\beta+u_i$, con $k$ parámetros:

\begin{center}
\begin{tabular}{@{}p{3.2cm}p{5.2cm}p{5.2cm}@{}}
\toprule
 & Error $u_i$ & Residuo $\hat u_i=y_i-x_i'\hat\beta$ \\
\midrule
Distancia a & La recta verdadera ($\beta$) & La recta estimada ($\hat\beta$) \\
¿Se observa? & No & Sí \\
¿Depende de nuestra elección? & No: lo fija la naturaleza & Sí \\
\bottomrule
\end{tabular}
\end{center}

Hay un residuo por observación, $n$ en total. MCO elige la recta que minimiza la
\emph{suma} de residuos al cuadrado (no cada residuo por separado, lo cual es imposible si
los puntos no están alineados):
\[
\hat\beta=\arg\min_\beta\sum_{i=1}^n(y_i-x_i'\beta)^2 .
\]
Derivando respecto de cada $\beta_j$ se obtienen $k$ condiciones de primer orden,
\[
\sum_{i=1}^n x_{ij}\hat u_i=0,\quad j=1,\dots,k\qquad\Longleftrightarrow\qquad X'\hat u=0 .
\]
Con constante, los residuos suman cero; con cada regresor, los residuos no están
correlacionados con él.

\begin{teorema}[Estimador insesgado de $\sigma^2$]\label{teo:s2reg}
Bajo $\E[uu'\mid X]=\sigma^2I$,
\[
s^2=\frac{1}{n-k}\sum_{i=1}^n\hat u_i^2\qquad\text{cumple}\qquad\E[s^2]=\sigma^2 .
\]
\end{teorema}
\begin{proof}
Sea $M=I-X(X'X)^{-1}X'$, simétrica e idempotente, con $\hat u=Mu$. Entonces
$\sum\hat u_i^2=u'Mu$. Como un escalar es igual a su traza y $\tr(AB)=\tr(BA)$,
\[
\E[u'Mu]=\E[\tr(Muu')]=\tr(M\,\sigma^2I)=\sigma^2\tr(M),
\]
y $\tr(M)=n-\tr\big((X'X)^{-1}X'X\big)=n-k$.
\end{proof}

\begin{intuicion}[el ajuste perfecto y el sobreajuste]
Los $n$ residuos deben cumplir $k$ ecuaciones, así que solo $n-k$ son libres. Lo que queremos
estimar es la varianza de los \emph{errores}; los residuos son versiones achicadas de ellos,
porque MCO ajustó $k$ parámetros para hacerlos pequeños. El $n-k$ corrige ese achicamiento.

\textbf{El caso $k=n$.} Hacer que la recta pase por todos los puntos es resolver $n$
ecuaciones con $k$ incógnitas. Si $k=n$, siempre hay solución exacta, \emph{sean cuales sean
los datos}: los residuos son cero y el ajuste es ``perfecto''. Pero
\[
\hat\beta=X^{-1}y=\beta+X^{-1}u :
\]
todo el ruido de la muestra quedó dentro de $\hat\beta$. En cambio, con $k<n$,
\[
\hat\beta-\beta=\Big(\tfrac1nX'X\Big)^{-1}\tfrac1n\sum x_iu_i ,
\]
y el último factor es un \emph{promedio}: los ruidos positivos y negativos se cancelan.

\textbf{Señal y ruido.} La señal ($x'\beta$) es común a todas las observaciones; el ruido es
propio de cada una. Solo comparando muchas observaciones se distingue lo que se repite (señal)
de lo que cambia (ruido), como con varios testigos de un mismo hecho. Con $k=n$ cada
parámetro tiene, en efecto, un solo testigo: el modelo memoriza señal y ruido mezclados.
Dentro de la muestra parece perfecto; con datos nuevos, cuyo ruido es otro, predice mal. Eso
es el \textbf{sobreajuste}. Los grados de libertad son los testigos de sobra que permiten
promediar y cancelar el ruido.
\end{intuicion}

\subsection{Posición, escala y forma: estandarizar}\label{sec:estandarizar}

Toda variable aleatoria se puede descomponer en tres piezas:
\[
Y=\underbrace{\mu}_{\text{posición}}+\underbrace{\sigma_Y}_{\text{escala}}\times
\underbrace{Z}_{\text{forma}},\qquad \E[Z]=0,\ \Var(Z)=1 .
\]
\begin{itemize}
  \item \textbf{Posición:} dónde está centrada.
  \item \textbf{Escala:} qué tan ancha es (la desviación estándar).
  \item \textbf{Forma:} lo que queda al quitar posición y escala; el patrón de cómo se reparte
        la probabilidad (campana, plana, asimétrica\ldots).
\end{itemize}
Cambiar la escala estira o comprime la distribución, pero no cambia su forma, igual que
cambiar el tamaño de una foto.

\begin{definicion}[Estandarizar]\label{def:estand}
\[
Z=\frac{Y-\mu}{\sigma_Y}.
\]
Restar la media quita la posición; dividir entre la desviación estándar quita la escala. Lo
que queda es la forma. Des-estandarizar es el camino inverso: $Y=\mu+\sigma_YZ$.
\end{definicion}

Aplicado a la media muestral, cuya desviación estándar es $\sigma/\sqrt n$:
\[
Z_n=\frac{\bar X_n-\mu}{\sigma/\sqrt n}=\frac{\sqrt n(\bar X_n-\mu)}{\sigma}.
\]
\textbf{De aquí viene el $\sqrt n$}: es parte de dividir entre la desviación estándar del
estimador.

\begin{intuicion}[el puente a la teoría asintótica]
Con cuentas de esperanza y varianza conocemos, para cualquier $n$, el \emph{centro} de
$\bar X_n$ ($\mu$) y su \emph{ancho} ($\sigma/\sqrt n$). Lo que en general no conocemos es su
\emph{forma}: solo la sabemos exactamente si los datos son normales (en cuyo caso el
estadístico $t$ con $s$ sigue una $t$ de Student con $n-1$ grados de libertad). La teoría
asintótica resuelve eso.
\end{intuicion}

% ============================================================
\section{Nociones de convergencia}\label{sec:conv}
% ============================================================

\subsection{Sucesiones de números}

\begin{definicion}[Límite de una sucesión]\label{def:limite}
$a_n\to L$ si para todo $\epsilon>0$ existe $N$ tal que $|a_n-L|<\epsilon$ para todo $n>N$.
\end{definicion}

$\epsilon$ es la tolerancia (qué tan cerca queremos estar) y $N$ es desde cuándo se cumple.
Por ejemplo, $a_n=\frac{n}{n+1}\to1$: como $|a_n-1|=\frac1{n+1}$, basta $n>\frac1\epsilon-1$.

\subsection{¿Por qué con variables aleatorias es distinto?}

$X_n$ es una función de $\omega$: para cada muestra posible hay una sucesión de números
distinta. Ya no basta la definición anterior; hay que decidir \emph{en qué sentido} $X_n$ se
acerca a su límite. Hay cuatro sentidos, de más a menos exigente en general.

\subsection{Los cuatro modos de convergencia}

\begin{definicion}[Convergencia casi segura]\label{def:cs}
$X_n\cas X$ si $\P\big(\{\omega: X_n(\omega)\to X(\omega)\}\big)=1$: para casi todo estado
del mundo, la sucesión de números converge.
\end{definicion}

\begin{definicion}[Convergencia en probabilidad]\label{def:cprob}
$X_n\cp X$ (o $\plim X_n=X$) si, para todo $\varepsilon>0$,
\[
\P\big(|X_n-X|>\varepsilon\big)\to0
\qquad\Longleftrightarrow\qquad
\P\big(|X_n-X|\le\varepsilon\big)\to1 ,
\]
donde $|\cdot|$ es la norma euclidiana.
\end{definicion}

\begin{intuicion}
Para cada $\varepsilon$ fijo, $p_n=\P(|X_n-X|>\varepsilon)$ es una \emph{sucesión de números},
y la definición pide que vaya a cero en el sentido usual. Hay dos tolerancias: $\varepsilon$
(qué tan cerca) y la del límite de $p_n$ (qué tan rara debe ser la excepción). No pide que
$X_n=X$; pide que equivocarse por más de $\varepsilon$ se vuelva improbable. La forma con
``$\le\varepsilon$'' (la que usa Hansen) es el evento complementario.

Para vectores, la convergencia en probabilidad equivale a la de cada componente, porque
$|v_j|\le\|v\|\le\sum_j|v_j|$.
\end{intuicion}

\begin{definicion}[Convergencia en $q$-ésimo momento]\label{def:cLq}
Para $q\ge1$, $X_n\cL{q}X$ si $\E\big[|X_n-X|^q\big]\to0$. Con $q=2$ se llama convergencia en
media cuadrática.
\end{definicion}

\begin{definicion}[Convergencia en distribución]\label{def:cdist}
$X_n\cd X$ si $\P(X_n\le u)\to\P(X\le u)$ en todo punto $u$ donde la CDF de $X$ es continua.
La distribución de $X$ se llama \textbf{distribución asintótica} o límite de $X_n$.
\end{definicion}

\begin{intuicion}[convergencia en distribución]
Solo pide que los \emph{histogramas} (las CDF) se parezcan, no que los valores estén cerca.
Si $Z\sim N(0,1)$ y $X_n=-Z$, entonces $X_n\cd Z$ trivialmente, aunque $|X_n-Z|=2|Z|$.

\textbf{¿Por qué solo en puntos de continuidad?} Sea $X_n=1/n$ (constante) y $X=0$. Para todo
$u\neq0$, $F_n(u)\to F(u)$; pero en $u=0$, $F_n(0)=0$ para todo $n$ y $F(0)=1$. Es absurdo
decir que $1/n$ no converge a $0$: el problema es solo el punto donde $F$ salta, donde toda la
masa está concentrada y un corrimiento mínimo la deja de un lado o del otro. Si el límite es
normal, su CDF no tiene saltos y la convergencia debe darse en todo $u$.

\textbf{Uso práctico:} para $n$ grande, $\P(X_n\le u)\approx F(u)$. Así se calculan p-valores e
intervalos.
\end{intuicion}

\subsection{Desigualdades fundamentales}

\begin{teorema}[Desigualdad de Markov]\label{teo:markov}
Para cualquier variable $X$ y $q,\varepsilon>0$,
\[
\P\big(|X|>\varepsilon\big)\le\frac{\E|X|^q}{\varepsilon^q}.
\]
\end{teorema}

\begin{intuicion}
Si el ingreso promedio es 1000, a lo más el 20\% puede ganar más de 5000: si fueran más, solo
ellos ya aportarían más de $0.2\times5000=1000$ al promedio. Si la media es chica, no puede
haber mucha probabilidad lejos del cero. Markov es el \textbf{puente entre momentos y
probabilidades}, y no requiere conocer la distribución.
\end{intuicion}

\begin{teorema}[Jensen]\label{teo:jensen}
Si $f$ es convexa y $\E|X|,\E|f(X)|<\infty$, entonces $f(\E[X])\le\E[f(X)]$. Si $f$ es cóncava,
la desigualdad se invierte.
\end{teorema}

En una curva con forma de U, el promedio de las alturas queda por encima de la altura en el
promedio. Con $f(x)=x^2$: $(\E X)^2\le\E X^2$, es decir, $\Var(X)\ge0$.

\begin{teorema}[Cauchy--Schwarz]\label{teo:cs}
Si $\E[u^2],\E[v^2]<\infty$, entonces $\E[uv]^2\le\E[u^2]\E[v^2]$, con igualdad solo si
$u=\alpha v$ con probabilidad 1.
\end{teorema}
\begin{proof}
Para todo $\lambda$, $0\le\E[(u-\lambda v)^2]=\E[u^2]-2\lambda\E[uv]+\lambda^2\E[v^2]$. Una
parábola que nunca es negativa tiene discriminante no positivo:
$4\E[uv]^2-4\E[u^2]\E[v^2]\le0$.
\end{proof}

Aplicada a variables centradas, dice que $|\Corr(X,Y)|\le1$: cuánto se mueven juntas dos
variables nunca supera el producto de sus tamaños.

\subsection{¿Qué es aleatorio y qué es el límite?}

Lo aleatorio es siempre $X_n$, lo que calculamos con la muestra. El nombre del modo
(``en probabilidad'', ``en distribución'') describe \emph{el criterio} con que juzgamos la
cercanía, no la naturaleza del límite. En econometría, sin embargo, hay un uso típico:

\begin{center}
\begin{tabular}{@{}llll@{}}
\toprule
Modo & Lo aleatorio & El límite & Para qué \\
\midrule
En probabilidad & $\hat\beta_n$ & $\beta$ (constante) & Consistencia \\
En distribución & $\sqrt n(\hat\beta_n-\beta)$ & $N(0,V)$ (aleatorio) & Inferencia \\
\bottomrule
\end{tabular}
\end{center}

Una constante es una variable aleatoria \textbf{degenerada}: toda su probabilidad está en un
punto, su varianza es cero y su CDF es un escalón. Es el caso extremo de una distribución a la
que se le quitó toda la dispersión.

% ============================================================
\section{Teoría asintótica: las herramientas}\label{sec:asin}
% ============================================================

\subsection{Ley débil de los grandes números}

\begin{teorema}[WLLN]\label{teo:wlln}
Si $X_1,\dots,X_n$ son i.i.d. y $\E[X_i]=\mu$ existe, entonces $\bar X_n\cp\mu$.
\end{teorema}
\begin{proof}[Demostración (con varianza finita)]
Por Markov con $q=2$ aplicado a $\bar X_n-\mu$ y el Teorema \ref{teo:mediamuestral},
$\P(|\bar X_n-\mu|>\varepsilon)\le\frac{\sigma^2}{n\varepsilon^2}\to0$. El enunciado general
solo requiere que exista la media.
\end{proof}

\begin{intuicion}
Al promediar muchas observaciones independientes, los errores por arriba y por abajo se
cancelan, y el pedazo (la muestra) se parece cada vez más al todo (la población). Es el
\textbf{motor} de casi todos los estimadores: el parámetro se escribe con promedios
poblacionales y el estimador usa los mismos promedios, pero muestrales.
\end{intuicion}

\subsection{Consistencia}

\begin{definicion}[Consistencia]\label{def:consist}
$\hat\theta_n$ es consistente para $\theta$ si $\hat\theta_n\cp\theta$, donde $\theta$ es el
valor \emph{verdadero} del parámetro.
\end{definicion}

\begin{intuicion}[el arquero]\leavevmode
\begin{itemize}
  \item \textbf{Insesgado:} los tiros, aunque dispersos, están centrados en el blanco
        ($\E[\hat\theta_n]=\theta$ para cada $n$). Es una propiedad de muestra finita.
  \item \textbf{Consistente:} con más práctica (más datos), los tiros se juntan sobre el blanco.
        Es una propiedad del límite.
\end{itemize}
Ninguna implica la otra: $X_1$ es insesgado pero no consistente; $\bar X_n+\frac1n$ es sesgado
pero consistente. Y converger no basta: $\bar X_n+1$ converge, pero a $\mu+1$, así que no es
consistente.
\end{intuicion}

\subsection{Convergencia conjunta y teorema de la función continua}

\begin{proposicion}[Marginal implica conjunta (en probabilidad)]\label{prop:conj}
Si $X_n\cp X$ e $Y_n\cp Y$, entonces $(X_n,Y_n)\cp(X,Y)$.
\end{proposicion}
\begin{proof}
$|(X_n,Y_n)-(X,Y)|\le|X_n-X|+|Y_n-Y|$. Si la suma supera $\varepsilon$, al menos uno de los
términos supera $\varepsilon/2$, así que
$\P(\cdot>\varepsilon)\le\P(|X_n-X|>\tfrac\varepsilon2)+\P(|Y_n-Y|>\tfrac\varepsilon2)\to0$.
\end{proof}

\begin{teorema}[Teorema de la función continua (CMT)]\label{teo:cmt}
Si $g$ es continua en un conjunto $C$ con $\P(X\in C)=1$, entonces
\[
X_n\cp X\Rightarrow g(X_n)\cp g(X),\qquad X_n\cd X\Rightarrow g(X_n)\cd g(X).
\]
\end{teorema}

\begin{intuicion}
Si los ingredientes llegan a su valor verdadero, una receta continua también llega al
resultado verdadero: un error pequeño en los ingredientes produce un error pequeño en el
resultado. Junto con la propiedad anterior, permite combinar: $X_n+Y_n$, $X_nY_n$,
$X_n/Y_n$ (si $Y\neq0$), $1/\bar X_n$, etc.

\textbf{¿Por qué continuidad?} Con $g(x)=\mathbf 1\{x>0\}$ y $X_n=1/n\to0$: $g(X_n)=1$ para todo
$n$, pero $g(0)=0$. El salto de $g$ justo en el límite rompe el resultado.

En distribución: si $Z_n\cd N(0,1)$, entonces $Z_n^2\cd\chi^2_1$ y $|Z_n|\cd|N(0,1)|$.
\end{intuicion}

\begin{proposicion}[Propiedades de la convergencia en distribución]\label{prop:propd}\leavevmode
\begin{enumerate}
  \item La convergencia marginal en distribución \emph{no} implica la conjunta. Con
        $Z\sim N(0,1)$: los pares $(Z,Z)$ y $(Z,-Z)$ tienen las mismas marginales, pero
        $Z+Z=2Z$ y $Z-Z=0$. La distribución de cada una no dice cómo se relacionan.
  \item Excepción: si $X_n\cd X$ e $Y_n\cd c$ (constante), entonces $(X_n,Y_n)\cd(X,c)$.
  \item Si $X_n\cd X$ y $|Y_n-X_n|\cp0$, entonces $Y_n\cd X$.
\end{enumerate}
\end{proposicion}

\subsection{El lema de Slutsky}

\begin{teorema}[Slutsky]\label{teo:slutsky}
Si $X_n\cd X$ e $Y_n\cp c$, con $c$ una constante, entonces
\[
X_n+Y_n\cd X+c,\qquad X_nY_n\cd cX,\qquad \frac{X_n}{Y_n}\cd\frac{X}{c}\ \ (c\neq0).
\]
Vale también con vectores y matrices de dimensiones compatibles: si $Y_n\cp C$, entonces
$Y_nX_n\cd CX$, $X_n'Y_n\cd X'C$ y $Y_n^{-1}X_n\cd C^{-1}X$ (si $C$ es invertible).
\end{teorema}

\begin{intuicion}[el problema que resuelve]
En econometría casi siempre combinamos una parte que sigue siendo aleatoria en el límite
(como $\sqrt n(\bar X_n-\mu)$, por el CLT) con una parte que converge a un número fijo (como
$s_n\to\sigma$). En general, conocer el límite de cada parte por separado \emph{no} basta para
conocer el de la combinación: el ejemplo $(Z,Z)$ frente a $(Z,-Z)$ muestra que el resultado
depende de cómo se relacionan.

Si una de las partes converge a una \textbf{constante}, el problema desaparece: una constante no
tiene aleatoriedad, así que no puede moverse junto con nada. Para $n$ grande, $Y_n$
prácticamente \emph{es} $c$. \textbf{Lo que converge a una constante se reemplaza por esa
constante.}
\end{intuicion}

\begin{observacion}[¿por qué $p$ para la constante y $d$ para la parte aleatoria?]\leavevmode
\begin{itemize}
  \item \textbf{Parte constante:} cuando el límite es constante, convergencia en probabilidad y
        en distribución son equivalentes: si el histograma colapsa en un pico sobre $c$, toda la masa
        está cerca de $c$. El requisito real es que
        el límite \emph{sea una constante}.
  \item \textbf{Parte aleatoria:} se pide lo mínimo, solo distribución, porque el CLT solo da eso:
        $\sqrt n(\bar X_n-\mu)$ no se pega a ninguna variable fija; solo su forma se estabiliza.
  \item \textbf{Resultado en distribución:} una cadena es tan fuerte como su eslabón más débil
        (con $Y_n=1$, $X_nY_n=X_n$). Y es justo lo que se necesita para inferencia: probabilidades
        del tipo $\P(|t|>1.96)$.
  \item Si $Y_n$ converge a algo \emph{aleatorio}, Slutsky no aplica.
\end{itemize}
\end{observacion}

\begin{center}
\begin{tabular}{@{}lll@{}}
\toprule
Aplicación & Parte aleatoria & Parte que va a una constante \\
\midrule
Estadístico $t$ & $\sqrt n(\bar X_n-\mu)$ & $1/s_n\to1/\sigma$ \\
Método delta (consistencia) & $\sqrt n(\hat\theta-\theta)$ & $1/\sqrt n\to0$ \\
Método delta (límite) & $\sqrt n(\hat\theta-\theta)$ & $g'(\theta^*)\to g'(\theta)$ \\
MCO & $\frac1{\sqrt n}\sum x_iu_i$ & $(\frac1nX'X)^{-1}\to Q^{-1}$ \\
Wald multivariado & $\sqrt n(\bar X_n-\mu)$ & $\hat\Sigma_n^{-1}\to\Sigma^{-1}$ \\
\bottomrule
\end{tabular}
\end{center}

\subsection{Dispositivo de Cramér--Wold}

\begin{teorema}[Cramér--Wold]\label{teo:cw}
$X_n\cd X$ si y solo si $t'X_n\cd t'X$ para todo vector $t$.
\end{teorema}

Un vector converge en distribución si y solo si toda combinación lineal de sus componentes
converge; es como conocer un objeto por todas sus sombras. Convierte problemas multivariados en
problemas univariados.

\subsection{Teorema central del límite}

\begin{teorema}[CLT]\label{teo:clt}
\textbf{Univariado.} Si $X_1,\dots,X_n$ son i.i.d. con media $\mu$ y varianza
$\sigma^2<\infty$,
\[
\sqrt n(\bar X_n-\mu)\cd N(0,\sigma^2),\qquad
\frac{\sqrt n(\bar X_n-\mu)}{\sigma}\cd N(0,1).
\]
\textbf{Multivariado.} Si son vectores i.i.d. con $\Sigma$ finita,
$\sqrt n(\bar X_n-\mu)\cd N(0,\Sigma)$.
\end{teorema}
\begin{proof}[Demostración (del caso multivariado a partir del univariado)]
Para $t$ fijo, los $t'X_i$ son escalares i.i.d. con varianza $t'\Sigma t$, así que
$\sqrt n(t'\bar X_n-t'\mu)\cd N(0,t'\Sigma t)$, que es la distribución de $t'Z$ con
$Z\sim N(0,\Sigma)$. Como vale para todo $t$, Cramér--Wold da el resultado.
\end{proof}

\subsubsection*{¿Por qué no mirar directamente $\bar X_n-\mu$?}

Porque colapsa. Por la WLLN, $\bar X_n-\mu\cd0$: su varianza $\sigma^2/n$ se va a cero y su
histograma se aplasta en un pico (una distribución degenerada). Eso solo repite la
consistencia y no dice cuánto nos equivocamos con nuestro $n$.

\begin{center}
\begin{tabular}{@{}lll@{}}
\toprule
Lo que se mira & Límite & Qué dice \\
\midrule
$\bar X_n$ & $\mu$ (un pico) & Que llega a la verdad \\
$\bar X_n-\mu$ & $0$ (un pico) & Lo mismo, dicho de otra forma \\
$\sqrt n(\bar X_n-\mu)$ & $N(0,\sigma^2)$ & Cómo es el error: forma y tamaño \\
\bottomrule
\end{tabular}
\end{center}

\subsubsection*{¿Por qué justo $\sqrt n$?}

El error típico de $\bar X_n$ es $\sigma/\sqrt n$, y eso viene de la cancelación parcial de
errores independientes (Sección \ref{sec:media}). Multiplicar por $\sqrt n$ compensa
exactamente el encogimiento:
$\text{desv. est.}\big(\sqrt n(\bar X_n-\mu)\big)=\sqrt n\cdot\sigma/\sqrt n=\sigma$, fija.
Con un factor menor el error todavía colapsa; con uno mayor, explota.

\begin{intuicion}[posición, escala y forma]
Con la descomposición de la Sección \ref{sec:estandarizar},
\[
\bar X_n-\mu=\underbrace{\frac{\sigma}{\sqrt n}}_{\text{escala}}\times
\underbrace{Z_n}_{\text{forma}} .
\]
Lo que ``se encoge'' es la escala. Multiplicar por $\sqrt n$ la neutraliza y deja visible la
forma, que no depende de $n$. El $\sqrt n$ solo fija la escala; la \textbf{forma normal la da
el CLT}. Media $0$ y varianza $\sigma^2$ no hacen normal a una variable; lo que la hace normal
es su forma de campana.
\end{intuicion}

\subsubsection*{¿Por qué la forma es normal?}

Una suma sale extrema solo si todos sus términos salen extremos en la misma dirección a la
vez, lo cual es muy improbable; en cambio, hay muchísimas combinaciones que la dejan cerca
del centro. Esa cuenta de combinaciones produce siempre la misma forma, la campana, sin
importar la forma de cada pieza. La condición $\sigma^2<\infty$ garantiza que ninguna
observación pueda dominar la suma; con colas demasiado pesadas (Cauchy), el CLT falla.

\subsubsection*{Restar exactamente $\mu$}

Restar $\mu$ no hace normal a la variable (una normal puede tener cualquier media), pero es
necesario para que haya límite. Si restamos $c\neq\mu$,
\[
\sqrt n(\bar X_n-c)=\underbrace{\sqrt n(\bar X_n-\mu)}_{\cd N(0,\sigma^2)}
+\underbrace{\sqrt n(\mu-c)}_{\to\pm\infty},
\]
y la campana se escapa. Esto es lo que hace funcionar una prueba de hipótesis: si la nula es
falsa, el estadístico explota y terminamos rechazando.

\subsubsection*{¿Qué significa ``asintóticamente normal''?}

\begin{observacion}
$\hat\theta_n$ \textbf{no} converge a una normal: converge a la constante $\theta$ (un límite
degenerado). Lo que converge a una normal es el error ampliado $\sqrt n(\hat\theta_n-\theta)$.
``$\hat\theta_n$ es asintóticamente normal'' es una abreviación de esto último.
\end{observacion}

Son dos fotos del mismo estimador: sin zoom vemos \emph{a dónde llega} (un punto); con zoom
$\sqrt n$ vemos \emph{cómo se equivoca en el camino} (una campana). Para nuestra muestra,
des-estandarizamos, devolviendo la escala de nuestro $n$:
\[
\hat\theta_n\approx\theta+\frac{\sigma}{\sqrt n}Z,\quad Z\sim N(0,1)
\qquad\Longleftrightarrow\qquad
\hat\theta_n\approx N\Big(\theta,\frac{\sigma^2}{n}\Big).
\]
Esto no es un límite, es una aproximación para un $n$ dado. Su ancho es positivo para todo $n$
finito, y ese ancho es el error estándar.

\begin{observacion}[lo que el CLT no dice]
Es una aproximación: con $n$ finito el estimador no es exactamente normal, y la aproximación
empeora con datos muy asimétricos o de colas pesadas. No dice nada si la varianza es infinita.
Asume i.i.d.: con datos dependientes hay versiones del CLT que requieren estacionariedad y
sumabilidad de las autocovarianzas.
\end{observacion}

\subsection{Método delta}

\textbf{El problema.} Conocemos la distribución asintótica de $\hat\theta$, pero nos interesa
una función de él: $\theta^2$, $1/\theta$, $\log\theta$, una elasticidad, un multiplicador de
largo plazo.

\begin{intuicion}
Cerca de $\theta$, cualquier función suave se ve como una recta de pendiente $g'(\theta)$:
\[
g(\hat\theta)-g(\theta)\approx g'(\theta)\,(\hat\theta-\theta).
\]
El error en $g$ es el error en $\theta$ multiplicado por la pendiente: si $g$ es empinada
amplifica la incertidumbre; si es plana, la reduce. Una normal multiplicada por un número sigue
siendo normal, y su varianza se multiplica por el cuadrado de ese número.
\end{intuicion}

\begin{teorema}[Método delta (escalar)]\label{teo:delta1}
Si $\sqrt n(\hat\theta-\theta)\cd N(0,\sigma^2)$ y $g$ es continuamente diferenciable cerca de
$\theta$, entonces
\[
\sqrt n\big(g(\hat\theta)-g(\theta)\big)\cd N\big(0,\ g'(\theta)^2\sigma^2\big).
\]
\end{teorema}

\begin{proof}
\begin{enumerate}
  \item \textbf{Valor medio.} Existe $\theta^*$ entre $\hat\theta$ y $\theta$ tal que
        $g(\hat\theta)-g(\theta)=g'(\theta^*)(\hat\theta-\theta)$. Es Taylor de orden cero con
        resto de Lagrange: en vez de dejar un resto, la pendiente se evalúa en un punto
        intermedio elegido para que la igualdad sea exacta.
  \item \textbf{Multiplicar por $\sqrt n$:}
        $\sqrt n\big(g(\hat\theta)-g(\theta)\big)=g'(\theta^*)\cdot\sqrt n(\hat\theta-\theta)$.
  \item \textbf{Consistencia (Slutsky, primer uso).}
        $\hat\theta-\theta=\frac1{\sqrt n}\cdot\sqrt n(\hat\theta-\theta)\cd0\cdot N=0$; como el
        límite es constante, $\hat\theta\cp\theta$.
  \item \textbf{El punto intermedio.} $\theta^*$ está atrapado entre $\hat\theta$ y $\theta$, así
        que $\theta^*\cp\theta$. Como $g'$ es continua, por el CMT $g'(\theta^*)\cp g'(\theta)$.
  \item \textbf{Slutsky, segundo uso.} $g'(\theta^*)\cp g'(\theta)$ (constante) y
        $\sqrt n(\hat\theta-\theta)\cd N(0,\sigma^2)$, así que el producto converge a
        $g'(\theta)N(0,\sigma^2)=N(0,g'(\theta)^2\sigma^2)$.
\end{enumerate}
\end{proof}

Con Taylor de primer orden la prueba es análoga: $g(\hat\theta)=g(\theta)+g'(\theta)
(\hat\theta-\theta)+R$, con $R(h)=o(|h|)$; se escribe
$\sqrt nR=\sqrt n|\hat\theta-\theta|\cdot\frac{R}{|\hat\theta-\theta|}$, donde el primer factor
converge en distribución y el segundo en probabilidad a cero, y Slutsky elimina el resto.

\begin{ejemplos}\leavevmode
\begin{itemize}
  \item $g(u)=u^2$: $\sqrt n(\bar X_n^2-\mu^2)\cd N(0,4\mu^2\sigma^2)$.
  \item $g(u)=1/u$: $\sqrt n(1/\bar X_n-1/\mu)\cd N(0,\sigma^2/\mu^4)$; el signo de la derivada
        desaparece al elevarla al cuadrado.
  \item $g(u)=\log u$: $\sqrt n(\log\bar X_n-\log\mu)\cd N(0,\sigma^2/\mu^2)$.
  \item En la práctica: $\text{se}\big(g(\hat\theta)\big)\approx|g'(\hat\theta)|\,
        \text{se}(\hat\theta)$.
\end{itemize}
\end{ejemplos}

\begin{observacion}[el caso degenerado $g'(\theta)=0$]
Si la pendiente es cero (un máximo o un mínimo), la varianza límite es $0$: la recta tangente es
horizontal y no ve el error. Hay que usar Taylor de segundo orden y multiplicar por $n$ en vez
de $\sqrt n$; el límite resulta una chi-cuadrado.
\end{observacion}

\begin{teorema}[Método delta (vectorial)]\label{teo:deltak}
Sean $\theta\in\R^k$ y $g:\R^k\to\R^q$ continuamente diferenciable cerca de $\theta$. Si
$\sqrt n(\hat\theta-\theta)\cd\xi\sim N(0,V)$, entonces
\[
\sqrt n\big(g(\hat\theta)-g(\theta)\big)\cd G'\xi\sim N(0,\ G'VG),\qquad
G=\frac{\partial g(\theta)'}{\partial\theta}\ (k\times q).
\]
Equivalentemente, con la matriz jacobiana $Dg=G'$ (de $q\times k$), la varianza es
$Dg\,V\,Dg'$.
\end{teorema}

\begin{intuicion}
Es el caso escalar con el gradiente en lugar de la derivada y la matriz $V$ en lugar de
$\sigma^2$. Con dos parámetros,
$\Var(a_1\xi_1+a_2\xi_2)=a_1^2V_{11}+a_2^2V_{22}+2a_1a_2V_{12}$, que en forma matricial es
$G'VG$. El término nuevo es la covarianza: si los estimadores están correlacionados, sus
errores se refuerzan o se compensan.
\end{intuicion}


\end{document}
